% resume6An26w.tex への挿入用。ヘッダーとページ番号はレジュメ本体で付ける。
% このディレクトリで実行：platex -interaction=nonstopmode -halt-on-error smalltest5_review.tex
% 続けて実行：dvipdfmx smalltest5_review.dvi
\documentclass[a4paper,10pt,dvipdfmx]{jsarticle}
\usepackage{../an26handout}

\renewcommand{\handout}[2]{%
  \def\handoutnumber{第#1回}%
  \begin{center}
    {\Large\normalfont\mdseries\color{ink}第#1回\quad #2}\par
    \vspace{2pt}{\color{accent}\rule{0.88\linewidth}{1.2pt}}
  \end{center}\vspace{2pt}}
\renewcommand{\topic}[1]{\par\vspace{5pt}%
  {\large\normalfont\mdseries\color{accent}#1}\par\vspace{1pt}}

\begin{document}
\pagestyle{empty}
\handout{5}{小テスト：広義積分}

学籍番号：\underline{\hspace{40mm}}\hfill
氏名：\underline{\hspace{40mm}}

次の広義積分を求めてください。極限を用いた計算過程も書いてください。

\topic{問題}
\begin{enumerate}[label=\color{accent}\arabic*.,leftmargin=2.4em,itemsep=9pt]
\item $\displaystyle\int_0^\infty e^{-x}\,dx$
\par\vspace{3mm}

\item $\displaystyle\iint_{(0,\infty)^2}e^{-x-y}\,dx\,dy$
\par\vspace{3mm}

\item $\displaystyle\iint_{(0,1)\times(1,\infty)}\frac{1}{y^2\sqrt{x}}\,dx\,dy$
\par\vspace{3mm}
\end{enumerate}

\par\vspace{4mm}
{\color{accent}\rule{\linewidth}{0.6pt}}
\topic{小テストの解答}
\begin{enumerate}[label=\color{accent}\arabic*.,leftmargin=2.4em,itemsep=9pt]
\item 広義積分の定義により、
$$
 \begin{aligned}
 \int_0^\infty e^{-x}\,dx
 &=\lim_{R\to\infty}\int_0^R e^{-x}\,dx\\
 &=\lim_{R\to\infty}\left[-e^{-x}\right]_0^R
 =\lim_{R\to\infty}(1-e^{-R})=1.
 \end{aligned}
$$

\item $R,S>0$ とし、$D_{R,S}=[0,R]\times[0,S]$ とおくと
$$
 \begin{aligned}
 \iint_{D_{R,S}}e^{-x-y}\,dx\,dy
 &=\left(\int_0^R e^{-x}\,dx\right)
   \left(\int_0^S e^{-y}\,dy\right)\\
 &=(1-e^{-R})(1-e^{-S}).
 \end{aligned}
$$
被積分関数は非負であり、$R,S\to\infty$ として領域を広げると
$$
 \iint_{(0,\infty)^2}e^{-x-y}\,dx\,dy
 =\lim_{\substack{R\to\infty\\S\to\infty}}
   (1-e^{-R})(1-e^{-S})=1.
$$

\item $D=(0,1)\times(1,\infty)$、$0<\varepsilon<1<R$ とする。
$D_{\varepsilon,R}=[\varepsilon,1]\times[1,R]$ 上では
$$
 \begin{aligned}
 \iint_{D_{\varepsilon,R}}\frac{1}{y^2\sqrt{x}}\,dx\,dy
 &=\left(\int_\varepsilon^1\frac{dx}{\sqrt{x}}\right)
   \left(\int_1^R\frac{dy}{y^2}\right)\\
 &=\left[2\sqrt{x}\right]_\varepsilon^1
   \left[-\frac1y\right]_1^R
 =2(1-\sqrt{\varepsilon})\left(1-\frac1R\right).
 \end{aligned}
$$
被積分関数は非負であり、境界は積分値に影響しない。
$\varepsilon\to0+$、$R\to\infty$ として領域を広げると
$$
 \iint_D\frac{1}{y^2\sqrt{x}}\,dx\,dy
 =\lim_{\substack{\varepsilon\to0+\\R\to\infty}}
  2(1-\sqrt{\varepsilon})\left(1-\frac1R\right)=2.
$$
\end{enumerate}
\end{document}
